\part{Formatting the Source Code}\label{sec:FormattingTheSourceCode}
As I said already in the introduction, coding conventions are primarily guidelines that should help making code better to maintain. Usually, this starts with making the code better to read.

In this regard, a source file is like any other text document: it benefits from having a clear structure. And that structure is supported by the formatting of the text.

An article may have sections with headlines in bold face, the paragraphs are separated by blank lines and the first lines of them are indented by 0.3~cm, and there are rules regarding line length, hyphenation, and how to treat widows and orphans.

This part of the document deals with the proper formatting of \Java source code files and their structure with the goal to create human readable \textit{documents}.

\begin{quote}
``[…], we want to establish the idea that a computer language is not just a way of getting a computer to perform operations but rather that it is a novel formal medium for expressing ideas
about methodology.'' \\
Herold Abelson et. al.: \textit{Structure and Interpretation of Computer Programs}
\autocite{Sussman:StructureAndInterpretationOfComputerPrograms}
\end{quote}

\subsubsection{Principles}
\begin{enumerate}[label=P\arabic*.]
	\item{The structure of the program text (the source code) must support the flow of reading.
	
	As said in the \ngref, bullet point~\ref{lst:ZoP:Readablity}: \bdq{Readability counts}.}
	
	\item{The formatting of the program text has to be consistent.
	
	Permanent changes in the way things are presented will confuse the reader.}
	
	\item{Navigation through the source file should be easy.
	
	When loaded inside an IDE, navigating inside a Java source file is easy: click on the name of the respective program element and the cursor is moved to the corresponding line in the file.
	
	But when using a plain text editor, the search function is often the only tool you have. A consistent structure with headlines and a defined order would help here, too.}
	
	\item{It does not harm when the source code also looks beautiful.
	
	Bullet point~\ref{lst:ZoP:BeautifulVsUgly} from the \ngref~ is also about the visual appearance of the source code. We can discuss how useful it is to use ASCII art in a source file, but not that it can contribute to the beauty of the program text.}
\end{enumerate}

\chapter{Lines and Spaces}

\section{Indentation}\label{sec:Indentation}
Indentation\index{indentation}, together with empty lines, is the most important tool to structure any source code. For the Python\index{Python} programming languange, the indentation is even part of the syntax\autocite{W3SCHOOLS:PythonIndentation}.

\subsection{Indentation Style}\label{sec:IndentationStyle}
The indentation\index{indentation} of code blocks, based on their context, increases the readability of the code.

Obviously, \bdq{readability} is a relative term, depending from one's habits. So some people are more familiar with the Kernighan-Ritchie (K\&R) indentation style\index{Kernighan-Ritchie style}\index{K\&R style|see{Kernighan-Ritchie style}}

\begin{lstlisting}
public final void myMethod (){
    if (flag) {
        …
    }
}   //  myMethod()
\end{lstlisting}

as it is also shown in the \ccjpl, other prefer to have the curly braces on a line of their own (sometimes referred to as GNU or BSD style\index{GNU style}\index{BSD style}\footnote{All, BSD, GNU and K\&R styles, will define more than just how to place the curly braces. In addition, all three are originally part of code conventions for programming in C and/or \mbox{C++} that cannot applied to Java without modifications. See chapter \tqfullvref{sec:OtherProgrammingLanguages} on this topic.}):

\begin{lstlisting}
public final void myMethod()
{
    if( flag )
    {
    	…
    }
}   //  myMethod()
\end{lstlisting}

Both styles will work fine, but they look different. So as a lot of other guidelines from a code conventions document, too, this is a matter of taste. But as conformity supports readability, a main purpose of these guidelines is to ensure that all source code looks similar, and we have to determine the style to use.

I prefer the so-called GNU or BSD style\index{indentation!style} and recommend it to you as well, because it makes it easier to detect the begin and the end of a code block than the Kernighan-Ritchie style. It demands to have the opening and closing curly braces on a line of their own, like below. This let the corresponding opening and closing braces appear on the same column:

\begin{lstlisting}
public final class MyClass
{
    public final void myMethod()
    {
        if( flag )
        {
            /* do something */
        }
        
        for( final var i = 0; i < max; ++i )
        {
            /* do something else */
        } 
    }   //  myMethod()
}   //  class MyClass
\end{lstlisting}
\subsection{Indentation Size}\label{sec:IndentationSize}
Again, at the very end, the size for the indentation is a matter of taste. But using four (4) spaces\index{indentation!base unit} has been established as a de~facto standard for the base unit of indentation in \Java source code, also because the original \bdq{Code Conventions for the \Java Programming Language}\autocite{SUN_CODE_CONVENTIONS:Indentation}index{Code Conventions for the Java Programming Language} by SUN\index{SUN} stated it that way. Therefore I recommend that you also use four spaces (\lstinline|' '|, ASCII~32 or 0x20) as the base unit of indentation.

When you write code \textit{snippets} for a publication (for an article, as part of a program documentation, or for an answer on StackOverflow\autocite{STACKOVERFLOW}~…), sometimes using a base indentation unit of two~spaces would be better, as it spares some room. But reformatting existing code for this use case is in most cases something that a text editor can do for you.
\subsection{Indentation Character}\label{sec:IndentationCharacter}
As said, you should use four \textit{blanks}\index{indentation!character}; do \textit{not} use tabs (\lstinline|'\t'|, ASCII~9 or 0x09)! Tabs are \textit{not} allowed! This is different from the \bdq{Code Conventions for the \Java Programming Language}\autocite{SUN_CODE_CONVENTIONS:Indentation}index{Code Conventions for the Java Programming Language} that do not explicitly specify whether to use blanks or tabs.

Usually both, the indentation character and the indentation base unit can be configured in your IDE; in Eclipse\index{Eclipse} this is part of the \textit{Formatter} definition (\verb#Window|Preferences|Java|Code Style|Formatter#).

To configure this in IntelliJ~IDEA\index{IntelliJ IDEA}, it is part of the Code~Style configuration for Java (\verb#File|Settings|Editor|Code Style|Java|Tabs and Indents#).



%==============================================================================
% End of file!!
%==============================================================================

\section{White Space}\label{sec:WhiteSpace}
This chapter covers how to use white space other than the indentation blanks for the formatting of your \Java source code.

\subsection{Blank Lines}\label{sec:BlankLines}
Blank lines improve readability by setting off sections of code that are logically related. These are the principles on how to use blank lines.

\subsubsection{Principles}
\begin{enumerate}[label=P\arabic*.]
\item{Each single line comment or block comment should be preceded by a blank line, unless it is not preceded by another comment}
\item{Use a blank line between the various parts of a class declaration}
\item{A blank line separates multiple class and interface definitions}
\item{Methods are also separated by single blank lines}
\item{Logical sections inside a method can be separated through a blank line to improve readability}
\item{Blank lines can be used to group code lines logically}
\item{More than one blank line should be avoided} 
\item{If it seems feasible for a particular purpose, feel free to use ASCII art to express the separation of parts in the source file}
\end{itemize}

This is again different from what the \bdq{Code Conventions for the \Java Programming Language}\autocite{SUN_CODE_CONVENTIONS:BlankLines}\index{Code Conventions for the Java Programming Language} stated regarding blank lines.

%==============================================================================
% End of file!!
%==============================================================================

\subsection{Round Brackets}\label{sec:Parenthesis}
Round Brackets (or Parentheses)\index{round brackets}\index{parentheses|see{round brackets}} are of course no whitespace characters – in the opposite, they are important syntax elements. But they trigger the use of whitespace, in particular blank spaces, and therefore, they are mentioned in this chapter.

We distinguish between \textit{arithmetical} round brackets and \textit{parameter} round brackets; the first is used in (mathematical or logical) expressions to order the terms, the latter is used with method or constructor calls or in lambdas.

The \bdq{Code Conventions for the \Java Programming Language}\autocite{SUN_CODE_CONVENTIONS:BlankSpaces}\index{Code Conventions for the Java Programming Language} does not make this distinction, and it handles them differently than described here (basically in the same way as described here as arithmetical round brackets).

\paragraph{Arithmetical Round Brackets}\index{arithmetical round brackets}
An opening arithmetical round bracket is always \textit{preceded} by a blank or another opening round bracket and \textit{never} followed by a blank. A closing arithmetical round bracket is \textit{never preceded} by a blank, but \textit{always followed} either by another closing round bracket, a blank, or a newline (or a semicolon, in case the statement ends).

This is a sample for arithmetical round brackets: \lstinline|((5 + 7) * 9 + 3) / 4|.

\paragraph{Parameter Round Brackets}\label{sec:ParameterParenthesis}\index{parameter round brackets}
An opening parameter round bracket is \textit{never preceded} by a blank and \textit{always followed} by one, while a closing parameter parenthesis is \textit{always preceded} by a blank. The only exception is for empty arguments lists where the opening parenthesis is immediately followed by the closing parenthesis.

Of course this is also valid for the declaration of methods and constructors.

Regarding to this rule, \lstinline|if|\index{if}, \lstinline|for|\index{for}, \lstinline|while|\index{while}, \lstinline|try|\index{try-with-resources} (with resources) and \lstinline|switch|\index{switch} are treated like they are functions.\footnote{And also the keywords \lstinline|catch|\index{try!catch} and \lstinline|synchronized|\index{synchronized}, but they do not allow complex terms inside their 'argument list'.}

Applying these rules can improve the readability of complex terms drastically, compared with the format suggested by the \ccjpl, as the examples below will illustrate.

\clearpage

\begin{tabular}{p{0.45\textwidth} p{0.45\textwidth}}
\hline
\lstinline[keepspaces=true]|final double d = sin( (a + b) / c );|
& 
\lstinline[keepspaces=true]|final double d = sin ((a + b) / c);|
\\
\hline
\lstinline[keepspaces=true]|System.out.println( "Finished!"  );|
&
\lstinline[keepspaces=true]|System.out.println ("Finished!");|
\\
\hline
\lstinline[keepspaces=true]|set( get() );|
&
\lstinline[keepspaces=true]|set (get ());|
\\
\hline
\lstinline[keepspaces=true]|System.out.println( sin( ((a + b) * (c + 4) + c)) / d ) );|
&
\lstinline[keepspaces=true]|System.out.println (sin (((a + b) * (c + 4) + c)) / d));|
\\
\hline
\lstinline[keepspaces=true]|if( a < b ) c = (d + e) / sin( f );|
&
\lstinline[keepspaces=true]|if (a < b) c = (d + e) / sin (f);|
\\
\hline
\lstinline[keepspaces=true]|while( ((line = reader.readLine()) != null) && (++i < capacity) ) processLine( line );|
&
\lstinline[keepspaces=true]|while (((line = reader.readLine ()) != null)~&& (++i < capacity)) processLine (line);|
\\
\hline
\end{tabular}

\keeplisting{3}
A special case are round brackets used with a cast:
\begin{lstlisting}
final Object o = "Text";
final var s = (String) o;
\end{lstlisting}

Here an opening round bracker is never followed by a blank, and a closing one is never preceded by a blank.

\subsubsection{Lambda Expressions}
According to the rules above, a lambda expression\index{lambda} with more than one argument should be written like this:
\keeplisting{3}
\begin{lstlisting}
final BiFunction f1 = ( a, b ) -> (a + b) * PI;
final BiFunction f2 = ( a, _ ) -> a * PI;
final BiFunction f3 = ( _, _ ) -> PI;
\end{lstlisting}

But it was found that it is easier to read without the blanks in the parameter list,  at least for single-letter parameter names or the unnamed variable\index{unnamed variable} (the underscore that is provided instead of the parameter):
\keeplisting{3}
\begin{lstlisting}
final BiFunction f1 = (a,b) -> (a + b) * PI;
final BiFunction f2 = (a,_) -> a * PI;
final BiFunction f3 = (_,_) -> PI;
\end{lstlisting}

So I suggest this style for lambda expressions.

%==============================================================================
% End of file!!
%==============================================================================

\subsection{Blank Spaces}\label{sec:BlankSpaces}
The principles below are in place regarding the use of blank spaces, above and beyond to what was said already.

\subsubsection{Principles}
\begin{enumerate}[label=P\arabic*.]
\item{For a keyword like \lstinline|if|, \lstinline|while|, \lstinline|for| etc., followed by a opening round bracket, there is no blank space between the keyword and the round bracket, but  between that bracket and the argument, and another is between the argument and the closing round bracket. 

As a short it can be said that the round brackets for these keywords will be treated like parameter round brackets, as already stated in \tqfullref{sec:ParameterParenthesis}.\footnote{This is in opposite to the \ccjpl; there it is recommended to have a blank between the keyword and the opening round bracket while not using it between the method name and the opening round bracket here, just to distinguish between keywords and methods calls. I think that it is more important to distinguish between round brackets in arithmetical expressions and those in method calls.}

Examples:
\begin{lstlisting} 
while( true ) { … }
\end{lstlisting}
\begin{lstlisting} 
for( int i = 0; i < 10; ++i ) { … }
\end{lstlisting}
\index{java.lang!Exception}\index{java.lang!Throwable!printStackTrace()}
\begin{lstlisting} 
try
{
    …
}
catch( Exception e )
{
    e.printStackTrace();
}
\end{lstlisting}
\begin{lstlisting} 
if( list.isEmpty() ) { … }
\end{lstlisting}}

\item{A blank space has to appear after commas in argument lists. The exception are generics where the comma is not followed by a blank space in case there is more than one type.

\keeplisting{7}
Examples:
\index{java.util!Map}\index{java.util!HashMap}
\begin{lstlisting}
Map<String,String> map1 = new HashMap<String,String>();

// using the diamond operator:
Map<String,String> map2 = new HashMap<>(); 

map.put( key, value );
\end{lstlisting}}

\item{All binary operators except “\verb#.#” has to be separated from their operands by spaces.

\keeplisting{6}
Example:
\begin{lstlisting}
a += c + d;

a = (a + b) / (c * d);

printSize( "size is " + foo + "\n" );
\end{lstlisting}}

\item{Blank spaces should never separate unary operators such as unary minus, increment (“\verb#++#”), and decrement (“\verb#--#”) from their operands.

\keeplisting{4}
Example:
\begin{lstlisting}
while( d-- > s++ )
{
    ++n;
}
\end{lstlisting}}

\item{The expressions in a \lstinline|for| statement has to be separated by blank spaces after the semicolon.

\keeplisting{4}
Example:
\begin{lstlisting}
for( expr1; expr2; expr3 )
{
    …
}
\end{lstlisting}

For an extended for loop, blank spaces has to be placed around the colon, too:
\begin{lstlisting}
for( String s : stringArray ) System.out.println( s );
\end{lstlisting}}

\item{For ternary statements, the question mark (“?”) and the colon (“:”) has to be surrounded by blank spaces:

\keeplisting{2}
Example:
\begin{lstlisting}
int signum = d < 0 ? -1 : d > 0 ? 1 : 0;
\end{lstlisting}

In this sample, some round bracket would be helpful to increase readability.}
\item{Casts should be followed by a blank space. 

\keeplisting{3}
Examples:
\begin{lstlisting}
myMethod( (byte) aNum, (Object) x);
myMethod( (int) (cp + 5), ((int) (i + 3)) + 1 );
\end{lstlisting}}
\end{enumerate}


%==============================================================================
% End of file!!
%==============================================================================


%==============================================================================
% End of file!!
%==============================================================================

\section{Line Length}\label{sec:LineLength}
Several sources (including the original \ccjpl) recommend to avoid code lines longer than 80~characters, because they were not handled well by many terminals and tools.

This had been true at the time those documents were written\footnote{The original \ccjpl were published by SUN\index{SUN} in 1999~…}, but today's  graphical screens are much wider (sometimes even 200 and more characters can be put into one line), and modern tools do not even care about line length, so we do not want to recommend a fixed line length for source code.\index{line length}

\textbf{Note:} Sample code for the use in documentations, articles or alike \textit{should have} a limited line length – generally not more than 70~characters.\index{line length!articles and documentation}

Different to lines with program code, \textit{comment lines} are limited to 80~characters – or more precise, they have to end at column~80\index{line length!comments}. All lines with “ASCII graphics” are extended to that column. This is pertaining to the beginning comments (chapter \tqfullvref{sec:BeginningComment}), the structuring comments (chapter \tqfullvref{sec:StructuringComments}) and most single-line comments (chapter \tqfullref{sec:SingleLineComments}).

\subsection{Wrapping Lines}\label{sec:WrappingLines}
As we have an unlimited line length, it will not happen that an expression does not fit into a single line. Nevertheless you might consider to break long lines to improve the readability of the code. In this case, you should follow these general principles:

\subsubsection{Principles}
\begin{enumerate}[label=P\arabic*.]
\item{Only one statement per line!}
\item{Break after a comma.}
\item{Break before an operator.}
\item{Break before a dot.}
\item{Break after an opening parenthesis and before a closing one.}
\item{Prefer higher-level breaks to lower-level breaks.}
\item{Align the new line with the beginning of the expression at the same level on the previous line.}
\item{If the above rules lead to confusing code or to code that's squished up against the right margin, consider to break up the code into multiple parts, by introducing additional temporary variables.}
\end{enumerate}

\subsubsection*{Line Wrapping for Method Calls}
\keeplisting{4}
Here are some examples on how to break long method calls:
\begin{lstlisting}
//  USUAL
final var result = myMethod( longExpression1, longExpression2, 
    longExpression3, longExpression4, longExpression5 );
\end{lstlisting}
Using a single line for each argument makes especially sense in cases where one or more expressions are really long, or if a comment should be added to the argument:
\keeplisting{8}
\begin{lstlisting}
myProcedure
( 
    longExpression1,
    longExpression2,
    longExpression3,
    longExpression4, // some comment
    longExpression5
);
\end{lstlisting}

\keeplisting{9}
With a return value:
\begin{lstlisting}
final var result = myFunction
    ( 
        longExpression1,
        longExpression2,
        longExpression3,
        longExpression4, // some comment
        longExpression5
    );
\end{lstlisting}

\keeplisting{11}
When another method call is used as an argument\footnote{Usually, you should avoid to call another method in an arguments list, although sometimes this cannot be avoided (refer to \tqfullvref{sec:CodingConstructors} and there to section \tqfullref{sec:ConstructorHelpers}). But in most cases, introducing a local variable is the better option.}:
\keeplisting{9}
\begin{lstlisting}
final var result = myMethod( longExpression1,
    someMethod( longExpression2,
        longExpression3 ) );

final var result = myMethod1( longExpression1,
    someMethod( longExpression2,
        longExpression3 ),
    longExpression4,
    otherMethod( longExpression5 ) );
\end{lstlisting}

Next there are two examples of breaking an arithmetic expression. The first is preferred, since the break occurs outside the parenthesized expression, which is at a higher level.
\keeplisting{10}
\begin{lstlisting}
// PREFERRED
final var longName1 = longName2 * (longName3 + longName4 - longName5)
    + 4 * longname6;
           
longName1 = longName2 * (longName3 + longName4 - longName5)
    + 4 * longname6;
           
// AVOID!!!           
longName1 = longName2 * (longName3 + longName4
                       - longName5) + 4 * longname6;
\end{lstlisting}

\keeplisting{15}
Then we have a few examples for the indentation on long method declarations. First the conventional case:
\begin{lstlisting}
//  CONVENTIONAL INDENTATION
public final void myMethod( final int anArg, final Object anotherArg,
    final String yetAnotherArg, final Object andStillAnother )
{
    …
}   //  myMethod()

private final synchronized void horkingLongMethodName( final int arg,
   final Object anotherArg, final String yetAnotherArg, 
   final Object andStillAnotherArg )
{
    …
}   //  horkingLongMethodName()
\end{lstlisting}

Next, a \lstinline|throws| clause that will be either appended after the closing round bracket on the same line, or it is placed on its own line with an indentation of four blanks; in that case, the consecutive lines with parameters are indented by 8~spaces.
\keeplisting{13}
\begin{lstlisting}[numbers=left]
public final void myMethod( final int anArg, final Object anotherArg,
    final String yetAnotherArg, final Object andStillAnother,
    final Instant andTheFinalArg ) throws IllegalArgumentException
{
    …
}   //  myMethod()

// ADDITIONAL INDENTATION BECAUSE OF THE throws CLAUSE
public final void myMethod( final int anArg, final Object anotherArg,
        final String yetAnotherArg, final Object andStillAnother )
    throws IllegalArgumentException, IOException
{
    …
}   //  myMethod()
\end{lstlisting}
Usually, the second variant (starting at line~9) is better readable and should be preferred.

\keeplisting{11}
The final examples uses a single line for each argument:
\begin{lstlisting}[numbers=left]
private static final synchronized anotherHorkingLongMethodLine
(
    final int anArg,
    final Object anotherArg,  // a comment …
    final String yetAnotherArg,
    final Object andStillAnother
)
{
    …
}   //  anotherHorkingLongMethodLine()
\end{lstlisting}

\keeplisting{11}
Or again with a \lstinline|throws| clause:
\begin{lstlisting}[numbers=left]
private static synchronized anotherHorkingLongMethodLineThrows
(
    int anArg,
    Object anotherArg,
    @MyAnnotation String yetAnotherArg,
    Object andStillAnother
) throws IllegalArgumentException, IOException
{
    …
}   //  anotherHorkingLongMethodLineThrows()
\end{lstlisting}

For a method declaration, a comment for a formal parameter (as in line~4 of the first example) is usually obsolete, as the description for the parameters is a mandatory part of the documentation comments (see chapter \tqfullvref{sec:MethodComment} for the details).

But you should think about having each formal argument in a line of it own if you have annotations\index{annotation} on parameters, as shown in line~5 of the second example. Especially when the annotation\index{annotation} will have its own arguments.

\keeplisting{5}
Here are the acceptable ways to format ternary expressions\footnote{Read more about the ternary operator \lstinline|?| in chapter \tqfullvref{sec:TheTernaryOperator}.}:
\begin{lstlisting}[numbers=left]
final var alpha = (aLongBooleanExpression) ? beta : gamma;
                                           
final var alpha = (aLongBooleanExpression)
    ? beta
    : gamma;
\end{lstlisting}

And finally a bad example; the formula below calculates the date for Easter sunday for a given year, based on Gauss' Easter algorithm \autocite{WIKIPEDIA:DateOfEaster,WIKIPEDIA:Gaussche_Osterformel}.

\index{Gauss@Gauß, Carl Friedrich}\index{Lichtenberg, Heiner}\index{java.lang!IllegalArgumentException}\index{java.time!LocalDate}\index{java.time!LocalDate!minusDays()}\index{java.time!LocalDate!of()}\index{java.time!LocalDate!plusDays()}\index{java.time!Month!MARCH}\index{java.time!Year}\index{java.time!Year!getValue()}\index{java.time!Year!isAfter()}\index{java.time!Year!isBefore()}\index{org.tquadrat.foundation.lang!Objects}\index{org.tquadrat.foundation.lang!Objects!requireNonNullArgument()}
\begin{lstlisting}[numbers=left]
/**
 *  <p>{@summary This method calculates the date of Easter Sunday 
 *  for the given year.} The year has to be in the range from 1583 
 *  to 3900 (included).</p>
 *  <p>The resulting date is for the Gregorian calendar.</p>
 *  <p>The algorithm itself is not the original one published by Carl
 *  Friedrich Gauß first in 1800 (corrected version in 1816), but
 *  that one published by Heiner Lichtenberg in 1997.
 *
 *  @thanks Carl Friedrich Gauß
 *  @thanks Heiner Lichtenberg  
 *
 *  @param  year    The year for which the date of Easter Sunday is 
 *      wanted for.
 *  @return The date of Easter Sunday in the given year.
 *  
 *  @see <a href="https://de.wikipedia.org/wiki/Gau%C3%9Fsche_Osterformel">Gaußsche Osterformel</a>
 */
public static final LocalDate calcEasterDate( final Year year )
{
    //---* Check the arguments *-------------------------------------    
    if( requireNonNullArgument( year, "year" ).isBefore( Year.of( 1583 ) ) )
    {
        throw new IllegalArgumentException( "This method will work only for years greater than or equal to 1583" );
    }
    if( year.isAfter( Year.of( 3900 ) ) )
    {
        throw new IllegalArgumentException( "This method will work only for years less than or equal to 3900" );
    }

    /*
     * Initialise the return value with the last day of February and 
     * add the calculated number of days.
     */
    final var retValue = LocalDate.of( year.getValue(), MARCH, 1 )
        .minusDays( 1 )
        .plusDays( (21 + ((19 * (year.getValue() % 19) + (15 + (3 * (year.getValue() / 100) + 3) / 4 - (8 * (year.getValue() / 100) + 13) / 25)) % 30) - (((19 * (year.getValue() % 19) + (15 + (3 * (year.getValue() / 100) + 3) / 4 - (8 * (year.getValue() / 100) + 13) / 25)) % 30) + (year.getValue() % 19) / 11) / 29) + (7 - ((21 + ((19 * (year.getValue() % 19) + (15 + (3 * (year.getValue() / 100) + 3) / 4 - (8 * (year.getValue() / 100) + 13) / 25)) % 30) - (((19 * (year.getValue() % 19) + (15 + (3 * (year.getValue() / 100) + 3) / 4 - (8 * (year.getValue() / 100) + 13) / 25)) % 30) + (year.getValue() % 19) / 11) / 29) - (7 - (year.getValue() + year.getValue() / 4 + (2 - (3 * (year.getValue() / 100) + 3) % 4)) % 7)) % 7) );

    //---* Done *----------------------------------------------------
    return retValue;
}   //  calcEasterDate()
\end{lstlisting}

Even proper breaking of the term in line~37 will not help much to understand what's going on there. A much better formatted (and therefore more readable) version of the method \lstinline|calcEasterDate()| can be found in the listing \tqvref{listing:GaussEaster}, showing what is meant with \bdq{breaking up the code into multiple parts by introducing additional temporary variables}.

%==============================================================================
% End of file!!
%==============================================================================


%==============================================================================
% End of file!!
%==============================================================================


%==============================================================================
% End of file!!
%==============================================================================

\chapter{Java Source File Structure}

Basically, what we refer to as a \bdq{Java Source File} is defined by the \Java programming language as a \bdq{Compilation Unit}\index{compilation unit} \autocite{ORACLE_DOC_LANGUAGE_SPECIFICATION:CompilationUnit}.

\bdq{The \Java Language Specification} distinguishes between \bdq{Ordinary Compilation Units}\index{compilation unit!ordinary}, \bdq{Compact Compilation Units}\index{compilation unit!compact} and \bdq{Modular Compilation Units}\index{compilation unit!modular}. Here, we are only interested in the \bdq{Ordinary Compilation Units}\index{compilation unit!ordinary}.

A \bdq{Modular Compilation Unit}\index{compilation unit!modular} (a \verb#module-info.java#\index{module-info.java} file) defines the module\index{module} structure of an application. The formatting of such a file is described in \tqfullvref{sec:ModuleDefinition} and modules in general are discussed in \tqfullvref{sec:Modularisation}.

\bdq{Compact Compilation Units}\index{compilation unit!compact} or \bdq{compact source files} were introduced with Java~25 and are mainly meant to support writing small scripts in Java, without the overhead of a full-flegded Java program. They should never be used in the context of a bigger Java program or application.

\section{File Encoding}
The \verb#javac# Java compiler expects UTF-8\index{UTF-8} \autocite{UNICODE:CoreSpec,WIKIPEDIA:UTF8} encoded source files. Obviously, ASCII\index{ASCII} \autocite{WIKIPEDIA:ASCII} encoded files will work, too, as well as most ISO-8859\index{ISO-8859} encodings, as long as only the ASCII characters are used.

A resource bundle or properties file\index{resource file}\index{properties file} \autocite{ORACLE_DOC_RESOURCEBUNDLE_CLASS} is expected to be encoded with ISO-8859-1\index{ISO-8859-1} \autocite{WIKIPEDIA:ISO88591}, although this can be changed by setting a system property. Refer to \autocite{ORACLE_DOC_PROPERTYRESOURCEBUNDLE_CLASS} for the details. 

The encodings of other source files depend on their format and usage.

%==============================================================================
% End of file!!
%==============================================================================

\section{Source File Size}
Above a certain number of lines, source code files become unwieldy, but it is difficult to specify a concrete value for the maximum size of a Java source file. First, generated source files tend to have lots of lines because it is often much easier to use boiler-plate code here than to generate methods and call them instead. And in particular when dynamically generated code is used\footnote{Code that is regenerated for each build.}, it would not even hurt the DRY principle: a change to the repeated code can be performed on one single place, the generator.

Second, classes that have lots of attributes usually end up with lots of lines, too: the formatting style and coding guidelines proposed in this document require a least 5~lines for each attribute, 6~lines for each getter method and 11~lines for each setter method, always including all comments and blank lines. This is about 22~lines for each attribute.  One hundred attributes seems to be a lot, but I have seen more than one real life application where database tables had more than that number of columns that each needed to be reflected as an attribute of the class that should reflect the table. I do not think that this is a good design anyway, but if your new Java class has to reflect the legacy data model, it has to have that number of attributes, bringing you to about 2200~lines of code without any methods other than the already mentioned getters and setters.

I found a limit of about 2000~lines to be feasible; if a class becomes larger than this number of lines, you should consider to move some of the code into other classes. This could be a utility class\index{utility class}, a parent class (probably abstract), or the functionality would be better placed with another entity in general. Also externalising some functionality by implementing the command pattern\autocite{Gamma:DesignPatterns} might be an idea. 

\subsubsection{Principles}
\begin{enumerate}[label=P\arabic*.]
\item{A Java source file should not exceed the limit of 2000~lines. As shorter it is, as better.}
\item{Generated source files may have significant more lines.}
\item{Other formatting rules may not be sacrificed only to reduced the number of lines of a source file.}
\end{enumerate}

%==============================================================================
% End of file!!
%==============================================================================

\section{The File Structure}

According to the Java Language Specification\autocite{ORACLE_DOC_LANGUAGE_SPECIFICATION:CompilationUnit} does each (ordinary) compilation unit\index{compilation unit} consist of three parts, each of which is optional:

\begin{itemize}
\item{A package declaration, giving the fully qualified name of the package to which the compilation unit belongs.

A compilation unit that has no package declaration is part of an unnamed package.}
\item{\lstinline|import| declarations that allow classes and interfaces from other packages, and static members of classes and interfaces, to be referred to by using their simple names.}
\item{Top level declarations of classes and interfaces.}
\end{itemize}

Based on that, I define a Java source code file having \textit{five} parts, in the given order:
\begin{enumerate}[nosep]
\item{\nameref{sec:BeginningComment}}
\item{\nameref{sec:PackageStatement}}
\item{\nameref{sec:ImportStatements}}
\item{\nameref{sec:ClassAndInterfaceDeclarations}}
\item{\nameref{sec:ClosingComment}}
\end{enumerate}

\subsection{Beginning Comment}\label{sec:BeginningComment}
A Java project will not consist from Java source files only; aside the already mentioned module definition files(see \tqref{sec:ModuleDefinition}\index{module-info.java}), the package documentation files\index{package-info.java} (refer to \tqfullvref{sec:PackageDocumentation}), there can be various build files, scripts, resource and property files, and even source files other programming languages.

All source files of a project – not only the Java sources – have to start with a comment that gives the copyright\index{copyright} notice and the license\index{license} information\footnote{This is different from the Sun coding conventions\autocite{SUN_CODE_CONVENTIONS:BeginningComments} that recommends to put also the class name and the date into the beginning comment. I omitted the name of the class because this is obvious from the file name in case of a \lstinline|public| class, and it would confuse the reader if the file contains more than one class. The date is considered obsolete.}, if possible at all: it should be obvious that it is sometimes difficult to add this kind of information to a binary resource like an image or a sound file.

For a Java file (a file with a name ending on \verb#*.java#), that beginning comment looks like this:
\begin{lstlisting}[numbers=left,caption={Beginning Comment for a Java file}]
/*
 * ==================================================================
 * Copyright © <Year> <Copyright Notice>
 * ==================================================================
 *
 * <License Notice>
 */
\end{lstlisting}

The same comment block can be used for Groovy\index{Groovy} source files, Kotlin\index{Kotlin} source files, Gradle\index{Gradle} build files (\lstinline|build.gradle|\index{build.gradle}) and generally for all file formats that uses the pattern \lstinline|/* … */| for comments.

A \textit{resource} file\index{resource file}\autocite{ORACLE_DOC_RESOURCEBUNDLE_CLASS} can have XML format or it may contain simple key-value pairs. In the latter form, comments are prepended with the hash symbol (``\verb|#|''), and the beginning comment has the following form: 
\begin{lstlisting}[numbers=left,language=,caption={Beginning Comment for a Resource or a Script file}]
#
# ===================================================================
# Copyright © <Year> <Copyright Notice>
# ===================================================================
#
# <License Notice>
#
\end{lstlisting}

Obviously, this can be used for all other file formats that use the same comment style, like Unix shell scripts\index{shell script}.

For a resource file in XML\index{XML} format\footnote{For an XML file, the beginning comments follows the initial \lstinline|<?xml …?>|, so it is not starting on line~1. That's similar for HTMl files.} (or any other XML file that is part of the sources), as well as for HTML\index{HTML} files, the beginning comment would look like this:
\begin{lstlisting}[numbers=left,firstnumber=5,language=XML,caption={Beginning Comment for XML and HTML files}]
<!--
=====================================================================
 Copyright © <Year> <Copyright Notice>
=====================================================================

 <License Notice>
-->
\end{lstlisting}

\TeX/\LaTeX\index{TeX}\index{LaTeX} uses the per cent sign (``\verb#%#'') to indicate a comment. So the beginning comment for a \TeX/\LaTeX\index{TeX}\index{LaTeX} source file would be like this:
\begin{lstlisting}[numbers=left,language=,caption={Beginning Comment Resource file}]
%
% ===================================================================
% Copyright © <Year> <Copyright Notice>
% ===================================================================
%
% <License Notice>
%
\end{lstlisting}

\verb#<Year>#, \verb#<Copyright Notice># and \verb#<License Notice># has to be replaced by the appropriate texts. The special character \bdq{\copyright} is used not only because is looks better; it tells you immediately whether your environment is capable to deal with UTF-8 encoded files.

The lines with the equals signs (“===”) will end at column~80; refer to chapter \tqfullvref{sec:LineLength}.\footnote{Keep in mind that in this document the line length is set to 70~columns; this means that you cannot just copy and paste the comment blocks into your development environment.}

\subsubsection{Principles}
\begin{enumerate}[label=P\arabic*.]
\item{Each file that is part of the sources for the program/application to build should have a beginning comment.}
\item{That beginning comment must contain the copyright notice with the year and the copytight holder.}
\item{It should contain a license notice (either the license itself or a reference to it).}
\item{No other information should be added to this comment block.}
\end{enumerate}

%==============================================================================
% End of file!!
%==============================================================================

\subsection{Package Statement}\label{sec:PackageStatement}
The first non-comment line of a Java source file is the \lstinline|package| statement\index{package}. According to the Java Language Specification\autocite{ORACLE_DOC_LANGUAGE_SPECIFICATION:UnnamedPackages}, it is optional (resulting in an \bsq{unnamed} or default package), but this is never used for production code. So you must have a \lstinline|package| statement in each and every Java source file.

For the names of packages, see chapter \tqfullvref{sec:Packages}.

\subsubsection{Principles}
\begin{enumerate}[label=P\arabic*.]
\item{Each Java source file has a \lstinline|package| statement.}
\end{enumerate}

%==============================================================================
% End of file!!
%==============================================================================

\subsection{Import Statements}\label{sec:ImportStatements}
Usually several \lstinline|import| statements will follow the \lstinline|package| statement. Only very basic code will not import other stuff than that from the \lstinline|java.lang|\index{java.lang} package. But interfaces with only one method (a so-called \bdq{functional interface}\index{functional interface} \autocite{ORACLE_DOC_FUNCTIONALINTERFACE_ANNOTATION}) may not need any imports at all, same for \bdq{marker interfaces} that does not declare any methods at all. Sometimes you may have also (abstract) parent classes without methods that do not need any \lstinline|import| statements, too.

\keeplisting{8}
But in general, that part of a source file would look like this:
\index{java.lang!String}\index{java.lang!String!format()}\index{java.lang!System}\index{java.lang!System!out}\index{java.awt.peer!CanvasPeer}\index{java.io!InputStream}
\begin{lstlisting}
import static java.lang.String.format;
import static java.lang.System.out;
…

import java.awt.peer.CanvasPeer;
import java.io.InputStream;
…
\end{lstlisting}
Static imports (for methods and constants) have to be placed before the imports for classes. Inside these blocks, the imports should be ordered alphabetically.

\keeplisting{6}
Wildcard imports like
\begin{lstlisting}
import static java.lang.String.*;
…

import java.util.*;
…
\end{lstlisting}
are not allowed and should not even be used for throwaway code. 

Eclipse users can use the function \verb#Source|Organize Imports# from the menu, or the short-key \verb#Shift+Ctrl+O#. This will reorder the import statements, explode the wildcards and remove unused imports; it will even add currently missing imports – given that this is properly configured in the \verb#preferences# at \verb#Java|Code Style|Organize Imports#.

At IntelliJ~IDEA, the function \verb#Code|Optimize Imports# from the menu does something similar; the short key would be \verb#Ctrl+Alt+O#.

\keeplisting{6}
Java~25 introduced a new feature, called \textit{Module Import Declarations} \autocite{JEP511}. It allows to import all public classes and interfaces of a module with a single statement:
\begin{lstlisting}
…
import module java.base;
…
\end{lstlisting}
As it can cause ambiguities, module imports should be avoided.

\subsubsection{Principles}
\begin{enumerate}[label=P\arabic*.]
\item{Place \lstinline|static| imports before non-static ones.}
\item{Sort the \lstinline|import| statements alphabetically.}
\item{Do never use wildcard imports.}
\item{Don't use module imports.}
\item{Remove unused imports. If an import is only required because of a Javadoc reference, change the Javadoc reference to use the fully qualified class name and remove the import.}
\end{enumerate}


%==============================================================================
% End of file!!
%==============================================================================

\subsection{Class and Interface Declarations}\label{sec:ClassAndInterfaceDeclarations}
Each Java source file\footnote{… or \bdq{ordinary compilation unit}\index{compilation unit}} contains at least one top level class or interface declaration; in most cases, it is just one. And it is strongly discouraged that it would be more than one!

Only one of these top level classes or interfaces can be \lstinline|public|, the others have to be package-local (meaning the modifiers \lstinline|private| and \lstinline|protected| are not allowed). If you really have more than one class in a single source file, the \lstinline|public| one should be the first entry, all others will follow, sorted alphabetically by their names.

With the JPMS\index{JPMS|see {Java Platform Module System}}\index{Java Platform Module System} (Java Platform Module System) the need for non-public top-level classes got nearly obsolete and they should be avoided.

If you need classes or interfaces that are closer related to a top-level class, you should consider to define these as \textit{inner} classes or interfaces. Inner classes can be \lstinline|public|, \lstinline|private| or \lstinline|protected|. An inner class is part of the top level class, but when defined as \lstinline|static|, it can be used independently from the enclosing class.

At least since Java~5, there are not only classes and interfaces that are declared in a Java source file. Today, we have

\begin{itemize}[nosep]
\item{\hyperref[subsec:Class]{classes} (keyword \lstinline|class|)}
\item{\hyperref[subsec:Interface]{interfaces} (keyword \lstinline|interface|)}
\item{\hyperref[subsec:Enum]{enums} (keyword \lstinline|enum|)}
\item{\hyperref[subsec:Record]{records} (keyword \lstinline|record|)}
\item{\hyperref[subsec:Annotation]{annotations} (keyword \lstinline|@interface|)}
\end{itemize}

Technically, records and enums are special types of classes, while an annotation\index{annotation} is somehow a special kind of an interface.

Each of these allows different sections in their declarations; the following tables describe the various sections and the order they should appear in. Each section has a header comment, and inside each section, the elements are ordered alphabetically.

\subsubsection{Class}\label{subsec:Class}
\keeplisting{34}
A skeleton for a Java \lstinline|class| may look like this:
\begin{lstlisting}[numbers=left,caption={Class Skeleton}]
public class MyClass 
{
        /*---------------*\
    ====** Inner Classes **==========================================
        \*---------------*/
    …    
        
        /*-----------*\
    ====** Constants **==============================================
        \*-----------*/
    …
        
        /*------------*\
    ====** Attributes **=============================================
        \*------------*/
    …
            
        /*------------------------*\
    ====** Static Initialisations **=================================
        \*------------------------*/
    …
        
        /*--------------*\
    ====** Constructors **===========================================
        \*--------------*/
    …
        
        /*---------*\
    ====** Methods **================================================
        \*---------*/
    …            
}
//  class MyClass
\end{lstlisting}
 
\begin{filecontents}{ClassParts.tbl}
  \begin{longtable}{|l|X|}
  \caption{Parts of a \lstinline|class| declaration} \\
  \hline 
  Part & Description \\ 
  \hline
  \endfirsthead
  \multicolumn{2}{c}%c
  {\tablename\ \thetable\ -- \textit{Continued from previous page}} \\
  \hline 
  Part & Description \\ 
  \hline
  \endhead
  \multicolumn{2}{r}{\textit{Continued on next page}} \\ 
  \endfoot
  \endlastfoot
  Inner Classes (optional) & Given the class defines inner classes, interfaces, enums or records, either \lstinline|private| or \lstinline|public|, they will come first. Internally, they will follow the same rules as for a top-level class. \\ 
  \hline 
  Constants (optional) & This part takes all constants defined by this class (meaning \lstinline|public static final| fields). All values for the constants are known at compile time, otherwise you should place the respective field to the \textit{Static Initialisations} part. \\ 
  \hline 
  Attributes (optional) & Here go the attributes for the class; technically, this part is in fact optional for a class, but usually a class without any attributes does not make much sense. Attributes are \lstinline|private| (under some exceptional circumstances, they can be \lstinline|protected|) and they can be \lstinline|static|, although this should be avoided. \\ 
  \hline 
  Static Initialisations (optional) & If a class declares \lstinline|static| fields that are not mere constants or their initialisation is more complex than just assigning a compile time constant, these fields go into this part.
  
  The fields should be initialised in a \lstinline|static {...}| block, not through an immediate initialiser.
  
  This is also the right place for the \lstinline|serialVersionUID|\index{serialVersionUID} of a seria­lisable\index{java.lang!Serializable} class. \\ 
  \hline 
  Constructors (optional) & All the constructors for the class go into this part. Technically, it is optional if the class will have only an empty \lstinline|public| default constructor, but the recommendation is to code this, too. \\ 
  \hline 
  Methods (optional) & All methods of the class are in this part. Again, this part is technically optional, but classes without methods are barely useful.  \\ 
  \hline 
 \end{longtable} 
\end{filecontents}
\LTXtable{\linewidth}{ClassParts.tbl}

\subsubsection{Interface}\label{subsec:Interface}
\keeplisting{20}
A skeleton for an \lstinline|interface| may look like this:

\begin{lstlisting}[numbers=left,caption={Interface Skeleton}]
public interface MyInterface 
{
        /*---------------*\
    ====** Inner Classes **==========================================
        \*---------------*/
    …
        
        /*-----------*\
    ====** Constants **==============================================
        \*-----------*/
    …
        
        /*---------*\
    ====** Methods **================================================
        \*---------*/
    …
                
}
//  interface MyInterface
\end{lstlisting}
 
\begin{filecontents}{InterfaceParts.tbl}
  \begin{longtable}{|l|X|}
  \caption{Parts of an \lstinline|interface| declaration} \\
  \hline 
  Part & Description \\ 
  \hline
  \endfirsthead
  \multicolumn{2}{c}%
  {\tablename\ \thetable\ -- \textit{Continued from previous page}} \\
  \hline 
  Part & Description \\ 
  \hline
  \endhead
  \multicolumn{2}{r}{\textit{Continued on next page}} \\ 
  \endfoot
  \endlastfoot
  Inner Classes (optional) & Defining inner classes for an interface is discouraged, although there are some valid use cases for this. So it could make sense to define an \lstinline|enum| with flag values used by a method declared by the interface. If the interface defines inner classes or interfaces, they have to be \lstinline|public| – classes have to be \lstinline|static|, too, and they will be the first part. Internally, they will follow the same rules as for top-level classes. \\ 
  \hline 
  Constants (optional) & This part takes all constants defined by this interface (meaning \lstinline|public static final| fields). All values for the constants are known at compile time. \\ 
  \hline 
  Methods (optional) & All methods of the interface are in this part. Again, this part is technically optional, it can be omitted for so-called \textit{marker interfaces}.  \\ 
  \hline 
 \end{longtable} 
\end{filecontents}
\LTXtable{\linewidth}{InterfaceParts.tbl}

\subsubsection{enum}\label{subsec:Enum}
\index{enum}An \lstinline|enum| is a special case of a Java class. Usually it only contains the enumeration values, but it can have more components. A skeleton for a \lstinline|enum| class may look like this:

\begin{lstlisting}[numbers=left,caption={enum Skeleton}]
public enum MyEnum 
{
        /*------------------*\
    ====** Enum Definitions **=======================================
        \*------------------*/
    …
        
        /*-----------*\
    ====** Constants **==============================================
        \*-----------*/
    …
        
        /*------------*\
    ====** Attributes **=============================================
        \*------------*/
    …
        
        /*------------------------*\
    ====** Static Initialisations **=================================
        \*------------------------*/
    …
        
        /*--------------*\
    ====** Constructors **===========================================
        \*--------------*/
    …
        
        /*---------*\
    ====** Methods **================================================
        \*---------*/
    …
                
}
//  enum MyEnum
\end{lstlisting}
 
\begin{filecontents}{EnumParts.tbl}
  \begin{longtable}{|l|X|}
  \caption{Parts of an \lstinline|enum| declaration} \\
  \hline 
  Part & Description \\ 
  \hline
  \endfirsthead
  \multicolumn{2}{c}%
  {\tablename\ \thetable\ -- \textit{Continued from previous page}} \\
  \hline 
  Part & Description \\ 
  \hline
  \endhead
  \multicolumn{2}{r}{\textit{Continued on next page}} \\ 
  \endfoot
  \endlastfoot
  Enum Definitions & This part is mandatory. Each \lstinline|enum| is an instance of the enum class. \\ 
  \hline 
  Constants (optional) & This part takes all additional constants defined by this \lstinline|enum| (meaning \lstinline|public static final| fields). All values for the constants are known at compile time, otherwise you should place the respective field to the \textit{Static Initialisations} part. \\ 
  \hline 
  Attributes (optional) & Here goes the attributes for the \lstinline|enum|, if any. Each attribute has to be \lstinline|private final| and will be initialised by the constructor. \\ 
  \hline 
  Static Initialisations (optional) & If an \lstinline|enum| declares \lstinline|static final| fields with an initialisation that is more complex than just assigning a compile time constant, these fields go into this part.
  
  The fields should be initialised in a \lstinline|static {...}| block. \\
  \hline 
  Constructors (optional) & All the constructors for the enum (rarely there is more than one) go into this part, if any. A constructor for an \lstinline|enum| must be \lstinline|private|. \\ 
  \hline 
  Methods (optional) & All methods of the \lstinline|enum| are in this part.  \\ 
  \hline 
 \end{longtable} 
\end{filecontents}
\LTXtable{\linewidth}{EnumParts.tbl}

Technically, an \lstinline|enum| can also define inner classes, but this is strongly discouraged, and I am not aware of any use case for this.

\subsubsection{Record}\label{subsec:Record}
\index{record}Records had been introduced to Java with version~14 and were finalised with version~16 \autocite{ORACLE_DOC_RECORD,ORACLE_DOC_LANGUAGE_SPECIFICATION:RecordClasses,Eckel:JavaRecords}. Like an \lstinline|enum|, a record is a restricted form of a Java class. It’s ideal for \bdq{plain data carrier} classes that contain data not meant to be altered. Usually it has only the most fundamental methods, such as constructors and accessors.

A skeleton for a Java record class may look like this, with all parts being optional:
\begin{lstlisting}[numbers=left,caption={Record Skeleton}]
public record MyRecord( ... ) 
{
        /*-----------*\
    ====** Constants **==============================================
        \*-----------*/
    …
        
        /*------------*\
    ====** Attributes **=============================================
        \*------------*/
    …
        
        /*------------------------*\
    ====** Static Initialisations **=================================
        \*------------------------*/
    …
        
        /*--------------*\
    ====** Constructors **===========================================
        \*--------------*/
    …
        
        /*---------*\
    ====** Methods **================================================
        \*---------*/
    …
                
}
//  record MyRecord
\end{lstlisting}
 
\begin{filecontents}{RecordParts.tbl}
  \begin{longtable}{|l|X|}
  \caption{Parts of a \lstinline|record| declaration} \\
  \hline 
  Part & Description \\ 
  \hline
  \endfirsthead
  \multicolumn{2}{c}%
  {\tablename\ \thetable\ -- \textit{Continued from previous page}} \\
  \hline 
  Part & Description \\ 
  \hline
  \endhead
  \multicolumn{2}{r}{\textit{Continued on next page}} \\ 
  \endfoot
  \endlastfoot
  Constants (optional) & This part takes all constants defined by this record (meaning \lstinline|public static final| fields). All values for the constants are known at compile time, otherwise you should place the respective field to the \textit{Static Initialisations} part. \\ 
  \hline 
  Attributes (optional) & Here go the attributes for the class; this part is optional, because usually, all attributes are defined as arguments to the \lstinline|record| definition. \\ 
  \hline 
  Static Initialisations (optional) & If a record declares \lstinline|static| fields that are not mere constants or their initialisation is more complex than just assigning a compile time constant, these fields go into this part.
  
  The fields should be initialised in a \lstinline|static {...}| block. \\ 
  \hline 
  Constructors (optional) & If the record needs additional constructors, these go into this part. It is optional, because in most cases no additional constructor is need. \\ 
  \hline 
  Methods (optional) & All additional methods of the record are in this part. \\ 
  \hline 
 \end{longtable} 
\end{filecontents}
\LTXtable{\linewidth}{RecordParts.tbl}

Same as \lstinline|enum|s can records have inner classes, but this is similarly discouraged, and there is no use case for it.

\subsubsection{Annotation}\label{subsec:Annotation}
Annotations had been introduced with Java~5 \autocite{ORACLE_DOC_LANGUAGE_SPECIFICATION:AnnotationInterfaces}. A skeleton for an annotation interface may look like this:

\keeplisting{8}
\begin{lstlisting}[numbers=left,caption={Annotation Skeleton}]
public @interface MyAnnotation 
{
        /*------------*\
    ====** Attributes **=============================================
        \*------------*/
    …
}
//  @interface MyAnnotation
\end{lstlisting}
 
\begin{filecontents}{AnnotationParts.tbl}
  \begin{longtable}{|l|X|}
  \caption{Parts of an annotation declaration} \\
  \hline 
  Part & Description \\ 
  \hline
  \endfirsthead
  \multicolumn{2}{c}%
  {\tablename\ \thetable\ -- \textit{Continued from previous page}} \\
  \hline 
  Part & Description \\ 
  \hline
  \endhead
  \multicolumn{2}{r}{\textit{Continued on next page}} \\ 
  \endfoot
  \endlastfoot
  Attributes (optional) & The attributes for the annotation go here, if any. An annotation without attributes would be a marker annotation. \\ 
  \hline 
 \end{longtable} 
\end{filecontents}
\LTXtable{\linewidth}{AnnotationParts.tbl}

\subsubsection{Principles}
\begin{enumerate}[label=P\arabic*.]
\item{Only one public top-level class (interface, record, enum, …) per compilation unit. Avoid non-public top-level classes.}
\item{Group the parts of the class as shown above. Remove the comments for unused sections.}
\end{enumerate}


%==============================================================================
% End of file!!
%==============================================================================


%==============================================================================
% End of file!!
%==============================================================================


%==============================================================================
% End of file!!
%==============================================================================

\chapter{Declarations}
This chapter describes how to format declarations of various kinds. See chapter \tqfullvref{sec:WrappingLines} for informations on how to break up the lines for long declarations.

\section{General Guidelines}

\subsection{Multiple Declarations per Line}
For local variables, it is possible to have multiple declarations in one single line, and the compiler does also not complain when you do this for attributes, too, but there it is even worse than for local variables.

So you should avoid the following:
\begin{lstlisting}
// AVOID!!!
int level, size;
\end{lstlisting}

It is strongly encouraged to have only one declaration per line, as it is also comforting to add some comments to the declaration. In other words, 
\begin{lstlisting}
int level; // indentation level
int size;  // size of table
\end{lstlisting}

is preferred over the sample above.

And by no means, you may put different types on the same line, despite the compiler does not complain about it. Example:
\begin{lstlisting}
// REALLY BAD!!!
int foo, fooarray [];
\end{lstlisting}

Adding comments to local variables should not be necessary if you name them appropriately (refer to \tqfullvref{sec:NamesForLocalVariables}). For attributes, the respective Javadoc comment (see \tqfullvref{sec:DocumentationComments}) is mandatory.

\subsubsection{Principles}
\begin{enumerate}[label=P\arabic*.]
\item{One declaration per line.}
\item{Comments for local variables are optional.}
\item{Documentation comments for attributes and constants are mandatory.}
\end{enumerate}

%==============================================================================
% End of file!!
%==============================================================================

\subsection{Placement of Declarations}

A declaration should be placed as close to the location where it is needed. This is mainly relevant for local variables, It means that you should wait with the declaration of variables until their first use, in particular when you need the result of an calculation for the initialisation. The \bdq{Code Conventions for the \Java Programming Language}\autocite{SUN_CODE_CONVENTIONS:Placement} recommends to put the declaration only to the beginning of a block.

For small blocks (this includes small method bodies), there is not much difference. Nevertheless, I find the Sun Code Conventions recommendation impractical, and I cannot follow the reasoning for that recommendation\footnote{To wait until first use \bdq{[…] can confuse the unwary programmer and hamper code portability within the scope}.}.

Some samples:
\begin{lstlisting}
for( var i = 0; i < limit; ++i )
{
    final var a = getValue( i );
    …
}
\end{lstlisting}
You declare the local variable even inside a loop; only the repeated \textit{initialisation} of a local variable inside a loop should be avoided, when possible:
\begin{lstlisting}[numbers=left]
// AVOID!
for( var i = 0; i < limit; ++i )
{
    final var a = new VeryComplexImmutableType();
    …
}
\end{lstlisting}
In this case, line~4 should be moved before the loop header (before line~2).


%==============================================================================
% End of file!!
%==============================================================================


%==============================================================================
% End of file!!
%==============================================================================

\section{Class Declarations}
When coding \Java classes, interfaces, enums, records, and annotations, you should follow the guidelines below:


\subsubsection{Principles}
\begin{enumerate}[label=P\arabic*.]
\item{.}
\item{.}
\item{.}
\end{enumerate}








\begin{itemize}
\item{No space between a method or constructor name and the opening parenthesis “(” starting its formal parameters list. Basically the parenthesis in the declaration are “paramater parenthesis” and should be treated along the same lines as for method and constructors calls; refer to chapter “\nameref{sec:ParameterParenthesis}” on page~\pageref{sec:ParameterParenthesis}.

Example:
\begin{lstlisting}
// WRONG!!
public abstract void myMethod (final int param);

// CORRECT
public abstract void myMethod( final int param );
\end{lstlisting}}

\item{If using generics, no space between the name of the declared class or interface and the opening “\textless”, and no blanks between the angle brackets, only around the keywords \lstinline|super| and \lstinline|extends|. Same when using the class or interface later.

Example:
\begin{lstlisting}
// WRONG
public interface MyInterface < T >
{ 
    … 
}
//  interface MyInterface

public OtherInterface < T, R > extends Function < T, R >, MyInterface < T >
{
    …
    
    public < U super T > U myMethod( final T arg );
}
//  interface OtherInterface

public MyClass implements otherInterface < Instant, String >
{
    …
}
//  class MyClass

// CORRECT
public interface MyInterface<T>
{ 
    … 
}
//  interface MyInterface

public OtherInterface<T,R> extends Function<T,R>, MyInterface<T>
{
    …
    
    public <U super T> U myMethod( final T arg );
}
//  interface OtherInterface

public MyClass implements otherInterface<Instant,String>
{
    …
}
//  class MyClass
\end{lstlisting}}

\item{An opening curly brace “\{” appears at a line by itself on the same indentation level as the declaration statement; anything afterwards will be indented with one additional level.

The corresponding closing curly brace “\}” starts a new line, indented to match its corresponding opening statement. For classes or interfaces, the name of the class or interface is repeated in a new line after the closing brace.

Example:\footnote{I omitted the structuring comments here; refer to chapters \tqfullvref{sec:ClassAndInterfaceDeclarations} and \tqfullvref{sec:StructuringComments}. And the bad sample is \textit{really} bad!}
\index{java.time!Instant!now()}\index{java.time!Instant!toString()}
\begin{lstlisting}
// AVOID!!
public final class MyClass {
public static InnerClass {
public final String innerMethod {
return Instant.now().toString(); } }
public static final void main( final String... args ) {
System.out.println( new InnerClass().innerMethod() ); } }

// RECOMMENDED
public final class MyClass 
{
    public static InnerClass 
    {
        public final String innerMethod 
        {
            return Instant.now().toString(); 
        }   //  innerMethod()
    }
    //  class InnerClass
    
    public static final void main( final String... args ) 
    {
        System.out.println( new InnerClass().innerMethod() ); 
    }   // main()
}
//  class MyClass    
\end{lstlisting}
Both versions compile properly, can be executed and will return the same result ~…}

\item{Inner classes, constants, attributes, constructors, methods, they all are separated from each other by a  single blank line (see chapter \tqfullref{sec:BlankLines}.}
\end{itemize}


%==============================================================================
% End of file!!
%==============================================================================


%==============================================================================
% End of file!!
%==============================================================================

\chapter{Statements}
This chapter deals with the formatting of the various types of statements.

\section{Simple Statements}
Each line may contain at most one single statement. Or, the other way round, each (simple) statement has to be written to its own line.

\keeplisting{7}
Example:
\begin{lstlisting}
// AVOID!!!
++argv; --argc;

// CORRECT
++argv;
--argc;
\end{lstlisting}

%==============================================================================
% End of file!!
%==============================================================================

\section{Compound Statements}\label{sec:CompoundStatements}
Compound statements\index{compound statement} are statements that contain lists of statements enclosed in curly braces \lstinline|{ <statement, … > }|. A compound statement (also called a \bdq{code block}\index{code block|see {compound statement}} or \bdq{block statement}\index{block statement| see {compound statement}}) typically appears as the body of another statement, such as the \lstinline|if| or \lstinline|for| statements, but it may appear also on it own\footnote{In Java, a block also indicates a scope\index{scope}; a local variable lives only in the scope it was declared in, and all inner scopes contained in it.}.

The block's opening curly brace \bdq{\{} appears alone at a line by itself with the same indentation as the statement it belongs to and the corresponding closing curly brace \bdq{\}} is placed at a new line with the same indentation. All code inside such a block is indented one level more than the curly braces.

Large compound statements (long code blocks) should be marked with a label\index{label} and should get an end comment, like in the examples below: 
\begin{lstlisting}
ForeverLoop: while( true )
{
    … // Lots of lines
}   //  ForeverLoop:

CountrySelector: switch( countryCode )
{
    case DE: …
    case EN: …
    … // All the other countries
    default: …
}   //  CountrySelector:

ProcessLoop: for( final var entry : entries )
{
    … // Lots of lines
}   //  ProcessLoop:
\end{lstlisting}

You may have seen something like this in existing code:
\begin{lstlisting}
// DISCOURAGED!
if( thisConditionIsTrue )
{
    // Lots of lines
    …
}   //	if( thisConditionIsTrue ) 
\end{lstlisting}
As long as the \lstinline|if| clause (or whatever statement starts the block) will not change, this pattern serves the same purpose as my variant with the label. But when you change the \lstinline|if|, you have to modify the end comment accordingly – something that is forgotten much too often!

As there is rarely a reason to change the label (the purpose of the block remains the same, no matter what the condition will be), it will not happen that the end comment will not match the start \bsq{comment} of the block.

See the chapters \tqfullref{sec:LabelsAndBreakStatements}, \tqfullref{sec:TrailingOrEndOfLineComments}, and \tqfullref{sec:CommentsWhen} for some additional details. 

For more examples regarding the formatting of compound statements refer in particular to the chapters \tqfullref{sec:IfStatements}, \tqfullref{sec:SwitchStatements}, \tqfullref{sec:ForStatements}, and \tqfullref{sec:WhileStatements}.

Empty compound statements must be avoided. At least a comment must be provided:
\begin{lstlisting}
{
    /* Nothing to do here */
}
\end{lstlisting}

%------------------------------------------------------------------------------

\subsection{\bsq{if}, \bsq{if-else}, \bsq{if-else~if-else} Statements}\label{sec:IfStatements}

We already discussed the use of blanks, round brackets, curly braces and linefeeds in \tqfullvref{sec:WhiteSpace}.

\keeplisting{11}
The base form for the \lstinline|if| statement is like this:
\begin{lstlisting}
if( <condition>)
{
    // <condition> is true
    <statements>;
}
else
{
    // <condition> is false
    <statements>;
}
\end{lstlisting}

\keeplisting{7}
If there is nothing to do in case the condition results as false, the \lstinline|else| clause must be omitted:
\begin{lstlisting}
if( <condition>)
{
    // <condition> is true
    <statements>;
}
\end{lstlisting}

\keeplisting{8}
Negate the condition if there is nothing to do when the original condition is true; do not let the first statement block empty:
\begin{lstlisting}
// RECOMMENDED!!
if( !<condition>)
{
    // <condition> is false; that means !<condition> is true
    <statements>;
}
\end{lstlisting}

\keeplisting{12}
Instead of:
\begin{lstlisting}
// AVOID!!
if( <condition>)
{
    // <condition> is true
    // Nothing to do here!
}
else
{
    // <condition> is false
    <statements>;
}
\end{lstlisting}

\keeplisting{3}
If there is no \lstinline|else| clause, and only a single statement, you can write the \lstinline|if| statement like this:
\begin{lstlisting}
if( <condition> ) <statements>;
\end{lstlisting}

The curly braces are mandatory if it is not possible to write the whole \lstinline|if| statement within a single line.

\keeplisting{22}
In case there are more conditions to check, but a \lstinline|switch| statement cannot be used due to the data types involved or the logic for the conditions, use this form:
\begin{lstlisting}
if( <condition1>)
{
    <statements>;
}
else if( <condition2>)
{
    <statements>;
}
else if( <condition3>)
{
    <statements>;
}
else if( <condition#>)
{
    <statements>;
}
else
{
    <statements>;
}
\end{lstlisting}

Of course you can repeat the \lstinline|else if| as often as you need, but at some point the readability suffers seriously, and you should think about another logic. Also the final \lstinline|else| clause is not mandatory.

\keeplisting{33}
It is not recommended to write \verb#if-else# like this, although it seems to be more conformant to the basic rule:
\begin{lstlisting}
// NOT RECOMMENDED!!
if( <condition1>)
{
    <statements>;
}
else 
{
    if( <condition2>)
    {
        <statements>;
    }
    else 
    {
        if( <condition3>)
        {
            <statements>;
        }
        else 
        {
            if( <condition#>)
            {
                <statements>;
            }
            else
            {
                <statements>;
            }
        }
    }
}        
\end{lstlisting}

The increasing indendation makes this even worse to read than the recommended variant, particularly when the code lines get longer.

\subsubsection{Principles}
\begin{enumerate}[label=P\arabic*.]
\item{Curly braces are mandatory when the \lstinline|if| statement has more than one line. They are always required for the \lstinline|else| clause.}
\item{No empty \lstinline|if| clause; negate the condition if necessary.}
\item{No blanks between the \lstinline|if|\index{if} keyword and the opening round bracket. 

Refer to \tqfullvref{sec:ParameterParenthesis} on how to place the white space.}
\item{The opening curly brace for the body statement of the \lstinline|if| and the optional \lstinline|else| clause is on the same indentation level as the keyword itself. In case of the \verb#if-else# construct, on the level of the \lstinline|else|.}
\item{The statements itself are indented one level relative to the curly braces.}
\end{enumerate}

Please refer also to chapter \tqvref{sec:TheTernaryOperator} that discusses the ternary \verb#?# operator.

%==============================================================================
% End of file!!
%==============================================================================

\subsection{\bsq{switch} Statements}\label{sec:SwitchStatements}
This chapter describes how to format a \lstinline|switch|\index{switch} statement together with the related \lstinline|case|\index{case} and \lstinline|default|\index{switch!default} statements.

Beginning with the first preview in Java~12, an alternative syntax for \lstinline|switch|\index{switch} was introduced, so that we have now two (in fact, three) different forms of that construct.

Some sources (e.g. \autocite{ORACLE_DOC_SWITCHEXPRESSIONS}) refer to the different syntax variants by the form of the \lstinline|switch|\index{switch} labels as ``\verb#case L:#'' (\bsq{old}) and ``\verb#case L ->#'' (\bsq{new}). 

\input{Import_SwitchTypesTable}

I will discuss in chapter \tqfullvref{sec:TheSwitchStatement} how and when to use which form.

%------------------------------------------------------------------------------

\input{formattingCode/SubSubSection_TraditionalSwitch}
\input{formattingCode/SubSubSection_NewSwitch}
\input{formattingCode/SubSubSection_PatternMatching}
\input{formattingCode/SubSubSection_SwitchFormatting}

%==============================================================================
% End of file!!
%==============================================================================

\subsection{‘for’ Statements}\label{sec:ForStatements}
Since Java~5 the language knows two distinct types of \lstinline|for|\index{for} loops, both described in \autocite{ORACLE_DOC_LANGUAGE_SPECIFICATION:For}.

The basic \lstinline|for| statement\index{for!basic} looks like this:
\begin{lstlisting}[numbers=left]
for( <initialization>; <condition>; <update> )
{
    <statements>;
}

for( <initialization>; <condition>; <update> ) <singleStatement>;
\end{lstlisting}
For the details refer to \autocite{ORACLE_DOC_LANGUAGE_SPECIFICATION:ForBasic}.

An enhanced \lstinline|for| statement\index{for!enhanced} (defined in \autocite{ORACLE_DOC_LANGUAGE_SPECIFICATION:ForEnhanced}) would be written like below:

\begin{lstlisting}[numbers=left]
for( <declaration> : <iterable> )
{
    <statements>;
}

for( <declaration> : <iterable> ) <singleStatement>;
\end{lstlisting}

The use of blanks, round brackets, curly braces and linefeeds was already discussed in \tqfullvref{sec:WhiteSpace}.

Curly braces can be omitted only when there is just a single simple statement to be executed in the loop. In addition, this single statement has to be written completely into the same line as the \lstinline|for| itself. In any other case, the curly braces are mandatory.

An empty loop body can be written like this:
\begin{lstlisting}[numbers=left]
for( <initialization>; <condition>; <update> )
{
    // Deliberately empty
}

for( <initialization>; <condition>; <update> );

for( <declaration> : <iterable> )
{
    // Deliberately empty
}

for( <declaration> : <iterable> );
\end{lstlisting}

In case the longer forms are used, the comment is mandatory, because empty code blocks are forbidden. I recommend the shorter forms (lines~6 and 13).

When explicitly using an iterator\index{iterator}\index{java.util!Iterator}\autocite{ORACLE_DOC_ITERATOR_INTERFACE} with the (basic) \lstinline|for|\index{for!basic} loop, this should look like this:
\begin{lstlisting}
for( final var i = source.iterator(); i.hasNext(); )
{
    …
}
\end{lstlisting}

When using the comma operator in the initialization or update clause of a \lstinline|for|\index{for} statement, avoid the complexity of using more than three variables. If needed, use separate statements before the \lstinline|for|-loop (for the initialization clause) or at the end of the loop (for the update clause).

Examples:
\begin{lstlisting}
for( var i = 0, j = target.length; i < source.length && j >= 0; ++i, --j )
{
    target [j] = source [i];
}

var proceed = true;
var i = source.iterator();
for( var j = 0; proceed; )
{
    target [j++] = i.next();
    proceed = i.hasNext() && j < target.length;
}
\end{lstlisting}

\subsubsection{Principles}
\begin{enumerate}[label=P\arabic*.]
\item{one}
\item{two}
\item{three}
\end{enumerate}

\lstinline|for| loops are discussed in more detail in \tqfullvref{sec:TheForStatement}.

Examples:
\begin{lstlisting}[numbers=left]
// Ok
for( var i = 0; i < max; ++i ) sum += i;
for( final var s : texts ) System.out.println( s );

// DISCOURAGED
for( var i = 0; i < max; ++i ) if( v [i] > 0 ) sum += v [i];
for( final var s : texts ) if( !s.empty() ) out.println( s );

// AVOID!!!
for( var i = 0; i < max; ++i ) if( v [i] > 0 )
{
    sum += v [i];
}
else
{
    sum -= v [i];
}

for( final var s : texts ) if( !s.empty() )
{
    System.out.print( "Value: " );
    System.out.println( s );
}
\end{lstlisting}
The samples in lines~6 and 7 are not correct according to this rule, because the \lstinline|if| statement is not a simple statement..

An empty \lstinline|for|-loop (one in which all the work is done in the initialization, condition, and update clauses) should have the following form:\footnote{An \textit{enhanced} \lstinline|for|-loop without body does not make that much sense (that mentioned above is a classical \lstinline|for|-loop). It may be possible to write an implementation of \lstinline|java.lang.Iterable| with an iterator that does something as a side effect of \lstinline|hasNext()| or \lstinline|next()|, but this would break the contract of the interface \lstinline|java.lang.Iterator|.}
\begin{lstlisting}
for( <initialization>; <condition>; <update> );
\end{lstlisting}


%==============================================================================
% End of file!!
%==============================================================================

\subsection{\bsq{while} and \bsq{do} Statements}\label{sec:WhileStatements}
The \lstinline|do|\index{while!do} statement does never come without a trailing \lstinline|while|\index{while} statement; this is defined in \autocite{ORACLE_DOC_LANGUAGE_SPECIFICATION:WhileDo}. The statement \lstinline|while|\index{while} itself is defined in \autocite{ORACLE_DOC_LANGUAGE_SPECIFICATION:While}. Both are loop statements, the difference is only where the termination condition is checked.

With \lstinline|do|\index{while!do} you define an exit-controlled loop\index{exit-controlled loop}\footnote{In German: \bdq{nicht-abweisende Schleife}.}; this means that the loop body is executed at least once:
\begin{lstlisting}
do
{
    <statements>;
}
while( <condition> );
\end{lstlisting}
Although writing
\begin{lstlisting}
do <statement>; while( <condition> );
\end{lstlisting}
is syntactically correct (note the semicolon after \verb#<statement># and before \lstinline|while|) you should always surround the loop body with curly braces when using \lstinline|do|\index{while!do}.

The entry-controlled \lstinline|while| loop\index{entry-controlled loop}\index{while}\footnote{In German: \bdq{abweisende Schleife}.} is used more often than the exit-controlled form; it looks like this:
\begin{lstlisting}
while( <condition> )
{
    <statements>;
}

while( <condition> ) <singleStatement>;
\end{lstlisting}

Same as for the \lstinline|for| statement (see \hyperref[sec:ForStatements]{there}), \bdq{single statement} also means \bdq{simple statement}. If the loop body is more complex, curly braces are mandatory; they are placed to lines of their own, indented to same level as the keyword, while the body itself has to be indented one additional level. The curly braces can be omitted only for a body consisting of a single statement that is written to the same line as the \lstinline|while|\index{while} statement.

The form
\begin{lstlisting}
// DISCOURAGED
while( <condition> )
    <statements>;
\end{lstlisting}
is syntactically correct, but have to be avoided. It is error prone in case the code will be adjusted at a later time.

The general use of blanks, round brackets, curly braces and linefeeds was already discussed in \tqfullvref{sec:WhiteSpace}.

When the loop body is larger, you should consider to use a label\index{label}, also when you use \lstinline|break|\index{break} or \lstinline|continue|\index{continue} inside the loop body:
\begin{lstlisting}[numbers=left]
<Label>: while( <condition> )
{
    if( <condition1> ) continue <Label>;
    if( <condition2> ) break <Label>;
    <statements>
}   //  <Label>:

<Label>: do
{
    if( <condition1> ) continue <Label>;
    if( <condition2> ) break <Label>;
    <statements>
} 
while( <condition> );
//  <label>
\end{lstlisting}

The body of a \lstinline|while| loop\index{while}\index{while!do} must not be empty! Or, more precisely, it must contain executable code! Usually a comment is sufficient to document that nothing happens inside that code block, but for a \lstinline|while|\index{while!empty loop} loop, this is not sufficient. The details are explained in the chapter \tqfullvref{sec:EmptyWhile}.

\subsubsection{Principles}
\begin{enumerate}[label=P\arabic*.]
\item{No blanks between the \lstinline|for|\index{for} and the opening round bracket. The blanks after the semi-colons are mandatory.

Refer to \tqfullvref{sec:ParameterParenthesis} on how to place the white space.}
\item{If the loop body is a simple single statement, it can be written without curly braces to the same line as the \lstinline|for|\index{for} itself.

Otherwise curly braces are mandatory.}
\item{The curly braces are on lines of their own, with the same indentation level as the \lstinline|for|\index{for} keyword. The code of the body will be indented one level.}
\item{Do not use a \lstinline|for|\index{for} loop to consume CPU cycles only!}
\item{Empty \lstinline|for|\index{for!empty} loops do all the work in the initialisation, condition and update clauses.

Prefer a \lstinline|while|\index{while} loop over an empty \lstinline|for|\index{for!empty} loop.}
\item{Keep the \lstinline|for|\index{for} statement simple.

Avoid using the comma operator in the initialisation and update clause.}
\item{Keep the loop body short.

Use a label\index{label} for longer loops.}
\item{Use the \lstinline|for|\index{for} loop variant that fits best to your current problem.

The enhanced \lstinline|for|\index{for!enhanced} is not generally better than the basic variant, it just matches better with some iteration tasks.}
\end{enumerate}

\lstinline|for| loops are discussed in more detail in chapter \tqvref{sec:TheForStatement}.










\subsection{‘do-while’ Statements}
A \lstinline|do-while| statement should look like below:







A \lstinline|while| statement should have one of the following forms:
\begin{lstlisting}[numbers=left]
while( <condition> )
{
    <statements>;
}
while( <condition> ) <singleStatement>
\end{lstlisting}

Again, the curly braces can only be omitted in case there is only a single statement, written on the same line than the \lstinline|while| statement that has to be executed in the loop.

An empty \lstinline|while|-loop\footnote{An empty \lstinline|while|-loop usually makes no sense as the compiler optimises it away. Only when the condition causes side effects, it will be executed. But such an implementation is difficult to understand and should be avoided.} should have the following form: 
\begin{lstlisting}
while( <condition> );
\end{lstlisting}

Longer \lstinline|while|-loops should use a label:
\begin{lstlisting}
LoopLabel: while( proceed )
{
    // Lots of code lines here …
}   //  LoopLabel:
\end{lstlisting}

Chapter \tqvref{sec:TheWhileStatement} has some more about \lstinline|while| loops.

%==============================================================================
% End of file!!
%==============================================================================

\subsection{\bsq{throw}, \bsq{try-catch} and \bsq{try-catch-finally} Statements}\label{sec:TheTryCatchStatements}
Throwing an exception and catching it somewhere else in the code is the foundation of error handling in the \Java programming language. The \verb#try-catch-finally#\index{try}\index{try!catch}\index{try!finally} statement is the core element here, together with the \lstinline|throw|\index{throw} statement and the exceptions itself. Error handling is covered in detail in chapter \tqfullvref{sec:ErrorHandling}.

The \lstinline|try|\index{try} statement is defined in \autocite{ORACLE_DOC_LANGUAGE_SPECIFICATION:Try}, and the definitions for the \lstinline|throw|\index{throw} statement can be found in \autocite{ORACLE_DOC_LANGUAGE_SPECIFICATION:Throw}.

\subsubsection{Formatting \bsq{throw}}
There is not much to say about the \lstinline|throw|\index{throw} statement in regards of formatting:
\begin{lstlisting}
throw <exception>;
\end{lstlisting}

That's it!

\subsubsection{Formatting \bsq{try-catch-finally}}
A \verb#try-catch-finally#\index{try}\index{try!catch}\index{try!finally} statement has the following format:
\begin{lstlisting}
try
{
    <statements>;
}
catch( final <ExceptionClass> e )
{
    <statements>;
}
finally
{
    <statements>
}
\end{lstlisting}
There can be more than one \lstinline|catch|\index{try!catch} clauses, but just one \lstinline|finally|\index{try!finally} clause. If there is at least one \lstinline|catch|\index{try!catch} clause, the \lstinline|finally|\index{try!finally} clause can be omitted. And \lstinline|catch|\index{try!catch} clauses are optional if there is a \lstinline|finally|\index{try!finally} clause.

A \verb#try-with-resources#\index{try-with-resources}\index{try} looks basically like this:
\begin{lstlisting}
try( <resource reference> )
{
    <statements>;
}
catch( final <ExceptionClass> e )
{
    <statements>;
}
finally
{
    <statements>
}
\end{lstlisting}
Again, there can be more than one \lstinline|catch|\index{try!catch} clauses, but just one \lstinline|finally|\index{try!finally} clause. But for a \verb#try-with-resources#\index{try-with-resources}\index{try}, both \lstinline|catch|\index{try!catch} clause and \lstinline|finally|\index{try!finally} clause can be omitted.

For a detailed description of \verb#try-with-resources#\index{try-with-resources}\index{try} refer to chapter \tqfullvref{sec:TryWithResources}.

The use of blanks, round brackets, curly braces and linefeeds was already discussed in \tqfullvref{sec:WhiteSpace}.

The curly braces around the code blocks for the \lstinline|try|\index{try}, the \lstinline|catch|\index{try!catch}, and the \lstinline|finally|\index{try!finally} are syntactically required. The opening and the closing braces will be placed to lines of their own, indented to the level of the respective keyword. The code inside will be indented one additional level, as for any other code block.

In case more than one exception type should be handled the same way, these can be combined into one \lstinline|catch|\index{try!catch} block, like this:
\begin{lstlisting}
try
{
    …
}
catch( final <ExceptionClass1> | <ExceptionClass2> e )
{
    <statements>;
}
\end{lstlisting}
The vertical bar is surrounded by blanks.

An \textit{empty} \lstinline|catch|\index{try!catch} block is unacceptable under all circumstances! Refer to chapter \tqfullvref{sec:GeneralExceptionHandling} how to handle exceptions. See also chapter \tqvref{sec:SingleLineComments} about empty blocks and comments.

Usually, the caught exception is referenced in the \lstinline|catch|\index{try!catch} block. If it is not referenced, you should use the underscore (\bdq{\lstinline|_|}) as the \textit{unnamed exception parameter}\autocite{ORACLE_DOC_LANGUAGE_SPECIFICATION:Declarations}\index{unnamed exception parameter}\footnote{From the \jls: \bdq{If a declaration does not include an identifier, but instead includes the keyword \bdq{\lstinline|_|} (underscore), then the entity cannot be referred to by name. […] A local variable, exception parameter, or lambda parameter that is declared using an underscore is called an unnamed local variable, unnamed exception parameter, or unnamed lambda parameter, respectively.}}. This looks like this:
\begin{lstlisting}
try
{
    …
}
catch( final <ExceptionClass> _ )
{
    <statements>;
}
\end{lstlisting}
Of course this is mandatory if the \lstinline|catch|\index{try!catch} block does not contain any executable code.

A \lstinline|finally| block without any executable code is completely obsolete and must be removed.

\subsubsection{Principles}
\begin{enumerate}[label=P\arabic*.]
\item{No blanks between the \lstinline|try|\index{try-with-resources} and the \lstinline|catch|\index{try!catch} keywords and the opening round brackets.

Refer to \tqfullvref{sec:ParameterParenthesis} on how to place the white space.}
\item{The keywords \lstinline|try|\index{try}, \lstinline|catch|\index{try!catch} and \lstinline|finally|\index{try!finally} are all indented at the same level.

The curly braces  for the respective code blocks are aligned with the keywords, the code is indented one level relative to the curly braces.}
\item{An empty \lstinline|catch|\index{try!catch} block is not allowed!

It must contain at least a comment that explains why the caught exception is not handled.}
\item{A \lstinline|finally|\index{try!finally} block without executable code must be removed.}
\end{enumerate}

%==============================================================================
% End of file!!
%==============================================================================

\subsection{‘synchronized’ Statements}\label{sec:SynchronizedStatements}
The \lstinline|synchronized|\index{synchronized} statement is defined in \autocite{ORACLE_DOC_LANGUAGE_SPECIFICATION:Synchronized}\footnote{The keyword \lstinline|synchronized|\index{synchronized} can be used also as a method modifier; this is explained in \autocite{ORACLE_DOC_LANGUAGE_SPECIFICATION:SynchronizedMethods}.}. A \lstinline|synchronized|\index{synchronized} statement acquires a mutual-exclusion lock on behalf of the executing thread, executes the related code block, then releases the lock. While the executing thread owns the lock, no other thread may acquire the lock – meaning no other thread can enter the block and execute the code in there.

The chapter \tqfullvref{sec:MultiThreading} discusses this topic further.

\keeplisting{5}
Such a \lstinline|synchronized|\index{synchronized} statement may look like this:
\begin{lstlisting}
synchronized( <expression> )
{
    <statements>
}
\end{lstlisting}
The curly braces are mandatory according the Java syntax; their placement follows the rules for any other code block. A single-line \lstinline|synchronized|\index{synchronized} statement is therefore not possible.

\keeplisting{8}
\jls \autocite{ORACLE_DOC_LANGUAGE_SPECIFICATION:Synchronized} says, that the type of \texttt{<expression>} has to be a reference type. That means that
\begin{lstlisting}
// DOES NOT WORK!!
final int i = 5;
synchronized( i )
{
    …
}
\end{lstlisting}
will not work\footnote{In fact, it will cause an error at compile time.}.

\keeplisting{6}
Unfortunately, the compiler will not complain about something like this:
\begin{lstlisting}
// DO NOT DO THIS!!
final StringBuilder sb = new StringBuilder();
synchronized( sb )
{
    …
}
\end{lstlisting}
\keeplisting{7}
And also not about this:
\begin{lstlisting}
// DO NOT DO THIS NEITHER!!
final String s = null;
synchronized( s )
{
    …
}
\end{lstlisting}
Ok, for the second one, you will get a \lstinline|NullPointerException|\index{java.lang!NullPointerException} at runtime, but it may take a lot of time before you spot the error in the first one. But at least most IDEs will issue respective warnings regarding both pattern (if not switched off~…). 

\keeplisting{7}
Even worse to debug would be this pattern, although it could be handy under some circumstances:
\begin{lstlisting}
// TRY TO AVOID!!
synchronized( provideLockObject() )
{
    …
}
\end{lstlisting}
In this case you have to rely on the method \lstinline|provideLockObject()| to return the proper guard object\index{guard object}.

So do not use local or temporary objects with the expression for \lstinline|synchronized|\index{synchronized}.

\keeplisting{16}
In case you declare a field only to provide the monitor for synchronisation, you should name it as \bsq{\texttt{Guard}}\index{guard object}.
\begin{lstlisting}
public final MyClass
{
    private final List<String> m_StringList = new ArrayList<>();
    private static final Object m_StringListGuard = new Object();
    
    public final void myMethod()
    {
        synchronized( m_StringListGuard )
        {
            …
        }
    }   //  myMethod()
}
//  class MyClass
\end{lstlisting}

\subsubsection{Principles}
\begin{enumerate}[label=P\arabic*.]
\item{No blanks between the \lstinline|synchronized|\index{synchronized} keyword and the opening round bracket.

Refer to \tqfullvref{sec:ParameterParenthesis} on how to place the white space.}
\item{The curly braces  for the code block are aligned with the keyword, the code is indented one level relative to the curly braces.}
\item{The \lstinline|synchronised|\index{synchronized} expression should be a final reference field, at best.

If that field exists only for this purpose, it should be named accordingly, as \bsq{\texttt{Guard}}\index{guard object}.}
\item{An empty \lstinline|synchronized|\index{synchronized} block does not make any sense and should be removed.}
\end{enumerate}

%==============================================================================
% End of file!!
%==============================================================================


%==============================================================================
% End of file!!
%==============================================================================

\section{Method Chaining}
Some APIs are designed to allow \bdq{method chaining}\index{method chaining}, sometimes also referred as \bdq{Call Chaining}\index{call chaining|see {method chaining}}; you may find this in particular for implementations of the \bdq{builder} pattern\index{builder pattern}, but it appears elsewhere, too. A prominent example is the class \lstinline|java.lang.StringBuilder|\index{java.lang!StringBuilder} \autocite{ORACLE_DOC_STRINGBUILDER_CLASS} (obviously not a builder, despite its name), a real example for a builder class would be the class \lstinline|DateTimeFormatterBuilder|\index{java.time.format!DateTimeFormatterBuilder} \autocite{ORACLE_DOC_DATETIMEFORMATTERBUILDER_CLASS} from the package \lstinline|java.time.format|\index{java.time.format}.

Usually, a class (or interface) will designed for method chaining\index{method chaining} when its methods return \lstinline|this| (the containing object instance), but other designs are possible that also allow method chaining\index{method chaining}, as shown in the code snippet below. It shows the use of method chaining\index{method chaining} on streams and on an instance of \lstinline|Optional|\index{java.util!Optional}.

\keeplisting{6}
\index{java.util!Map}\index{java.util!Map!Entry}\index{java.util!Optional}\index{java.util!Stream}
\begin{lstlisting}
return strings.stream()
    .collect( groupingBy( s -> s, counting() ) ) // returns java.util.Map
    .entrySet() 
    .stream()
    .max( Map.Entry.comparingByValue() ) // returns java.util.Optional
    .map( Map.Entry::getKey );
\end{lstlisting}

\keeplisting{3}
Method chaining allows you to write
\index{java.io!Writer}\index{java.io!Writer!append()}
\begin{lstlisting}
writer.append( "Caption\t: " )
    .append( value );
\end{lstlisting}

\keeplisting{3}
instead of
\index{java.io!Writer}\index{java.io!Writer!append()}
\begin{lstlisting}
writer.append( "Caption\t: " );
writer.append( value );
\end{lstlisting}

If the \bdq{chain} does not fit completely into a single line, each method call has to be written into a single line of its own.

\keeplisting{7}
In fact, it is recommended to do this all the time:
\index{java.io!Writer}\index{java.io!Writer!append()}
\begin{lstlisting}
// Acceptable
writer.append( "Caption\t: " ).append( value );

// RECOMMENDED
writer.append( "Caption\t: " )
    .append( value );
\end{lstlisting}

The dot has to be written in front of the method's name and the methods will be indented regularly.

\keeplisting{8}
Streams, that had been introduced with Java~8, will be formatted in the same way:
\index{java.util!Stream}
\begin{lstlisting}
final var names = customers.stream()
    .filter( c -> c.getCountry().equals( GERMANY )
    .filter( c -> c.getTurnover() > limit )
    .map( Customer::getName )
    .sorted()
    .toArray( String []::new );
\end{lstlisting}


\subsubsection{Principles}
\begin{enumerate}[label=P\arabic*.]
\item{Each method call after the first should be placed on a line of it own.

The dot prepends the method's name.}
\item{The method calls are indented one level relative to the first line of the call chain.} 
\end{enumerate}

%==============================================================================
% End of file!!
%==============================================================================

\section{‘return’ and ‘yield’ Statements}
The \lstinline|return|\index{return} statement has to be the last statement of a method, while the \lstinline|yield|\index{switch!yield} statement has to be the last statement of a \lstinline|case|\index{switch!case} block of a \lstinline|switch| expression\index{switch!expression}; refer to the chapters \tqfullvref{sec:ReturningValues} and \tqfullvref{sec:CaseResults} for more details on this.

When \lstinline|return|\index{return} is always the last statement of a method, it means that the simple form
\begin{lstlisting}
return;
\end{lstlisting}
is obsolete and will be omitted. This implies that there is only one \lstinline|return|\index{return} statement per method allowed, and that short-cutting a method is a forbidden practice.

Usually, the name of the return value is \lstinline|retValue| (see \tqfullvref{lbl:RetValue}), and the lines before the \lstinline|return|\index{return} statement are a blank line and the comment \lstinline|//---* Done *---…|:
\keeplisting{7}
\begin{lstlisting}
public final void myMethod()
{
    …
        
    //---* Done *----------------------------------------------------
    return retValue;
}   //  myMethod()
\end{lstlisting}

\lstinline|yield|\index{switch!yield} as a keyword was introduced with the new \lstinline|switch|\index{switch} format and \lstinline|switch| expressions\index{switch!expression} by Java~12. It is used to return a value from a \lstinline|case|\index{switch!case} block that is not an expression. Same as \lstinline|return|\index{return} has \lstinline|yield|\index{switch!yield}to be placed on a line of its own, a the end of a \lstinline|case| block\index{switch!case}. Without  a return value \lstinline|yield|\index{switch!yield} will always cause a compiler error. Some more details can be found at \autocite{ORACLE_DOC_LANGUAGE_SPECIFICATION:Yield}.

The use of round brackets with the \lstinline|return|\index{return} or \lstinline|yield|\index{switch!yield} statement is not required unless they make the return value more obvious and better readable in some way.

\subsubsection{Principles}
\begin{enumerate}[label=P\arabic*.]
\item{The \lstinline|return|\index{return} statement is, if it exists at all, always the last statement in a method.

There is at maximum one \lstinline|return|\index{return} statement in a message.}
\item{The \lstinline|return|\index{return} statement is introduced by a blank line and a \verb#Done!# comment.}
\item{A \lstinline|yield|\index{switch!yield} is always the last statement in a \lstinline|case|\index{switch!case} block.}
\item{The \lstinline|return|\index{return} and \lstinline|yield|\index{switch!yield} statements are always placed on a line of their own.}
\item{Do not use round brackets with \lstinline|return|\index{return} or \lstinline|yield|\index{switch!yield}.}
\end{enumerate}

%==============================================================================
% End of file!!
%==============================================================================


%==============================================================================
% End of file!!
%==============================================================================

\chapter{Embedded Code}\label{sec:EmbeddedCode}
Sometimes, it is necessary to embed code from another (programming) language into the \Java source code\footnote{If you write a code generator tool, this could even by Java code.}.

First, only relativly small fragments of foreign code should be embedded into the Java source. I consider something around 15 to 20~lines still as \bsq{small} in this context, but this can be discussed for sure.

Not only that the embedded code may clutter the Java source, it is also much better modifiable when stored in a resource file\index{resource file}. Therefore larger fragments and full documents can be handled better when provided as resources\index{resource file} (refer to \tqfullvref{sec:ResourcesClasspath})

\keeplisting{7}
Embedded code should be placed into a text block\index{text block}\footnote{For the definition of a text block\index{text block} in Java see \autocite{ORACLE_DOC_LANGUAGE_SPECIFICATION:TextBlocks}.}:
\begin{lstlisting}
final var code = """
````<line1>
````    <line2>
````        <line3>
````    <line4>
````""";
\end{lstlisting}

The left margin of the text inside the text block\index{text block} is indicated by the position of the closing three double quotes. In the sample above, \texttt{<line1>} in line~2 is not indented at all, while the lines~3 and 5 are indented by four spaces, and \texttt{<line3>} is indented by eight spaces. The indentation\index{indentation} within the source code does not matter here; it is just important for optical reasons. The \lstinline|" "| indicates the indentation blanks for the string only.

\keeplisting{5}
The embedded code looks like this when printed:
\begin{lstlisting}[language={}]
<line1>
    <line2>
        <line3>
    <line4>
\end{lstlisting}

Aligning the opening and the closing delimiters (the three double quotes) is not recommended. This would require reindentation of the text block\index{text block} if the variable name or modifiers are changed.

According to chapter \tqfullref{sec:IndentationCharacter} only blanks are used for the indentation\index{indentation!character} of the Java source code. But it is of course possible that the embedded code requires tabs for its indendation.
\keeplisting{6}
In this case, you have to the use the respective escape sequence \bdq{\lstinline|\\t|}:
\begin{lstlisting}[numbers=left]
final var codeJavaScript = """
````\tsetTimeout(function() {
````\t\tmyDisplayer("Hello!");
````\t}, 3000);
````""";
\end{lstlisting}

\bdq{Programmer's Guide to Text Blocks} \autocite{Laskey:TextBlocks} is a comprehensive guide to text blocks\index{text block} in \Java. Text blocks will also be discussed more generally again in chapter \tqfullvref{sec:TextBlocks}.

\subsubsection{Principles}
\begin{enumerate}[label=P\arabic*.]
\item{The opening three double-quotes (the \bsq{\textit{delimiter}}) for a text block\index{text block} should be placed at the right end of the previous line, and the closing delimiter should be placed on its own line, at the left margin of the text block.

Use the line-terminator escape sequence\footnote{A backslash immediately followed by a linefeed.} if the trailing line-feed is not wanted.}
\item{Avoid aligning the opening and closing delimiters and the text block's left margin.}
\item{If tabs are needed for the indentation of the code inside the text block, use the escape sequence \bdq{\lstinline|\\t|}. 

Mixing white space will lead to a result with irregular indentation.}
\item{Do not embed large code blocks.

Use resources for extensive code.}
\end{enumerate}

%------------------------------------------------------------------------------

\section{Formatting Embedding Code}\label{sec:FormattingEmbeddingCode}
If you have already defined formatting rules for the respective type of code you need to embed into your Java source code, you should first check if they can be applied to the embedded code, or if they can be modified to fit for this purpose. If so, use these rules and forget the following sections in this chapter.

If you have other types of code that need to be embedded on a regular basis, you should amend this document accordingly\footnote{Refer to \tqfullref{sec:HowToUseThisDocument}.}, with at least some basic rule on how to format the embedded code of that type.

%------------------------------------------------------------------------------

\subsection{Formatting SQL inside Java}\label{sec:FormattingSQLInsideJava}
If your code needs to access a database\index{database} on an RDBMS\index{RDBMS}, it is quite likely that you have to write some SQL\index{SQL} \autocite{ANSI:SQLStandard,W3SCHOOLS:SQL} statements into your code. Even when you are using an ORM\index{ORM} like JPA\index{JPA} \autocite{ECLIPSE_JPA} and Hibernate\index{Hibernate} \autocite{HIBERNATE_ORM}, it might still be necessary to add some SQL\index{SQL} statements or statements in the tool's variant of SQL (HQL\index{HQL}\index{HQL|see {Hibernate}} in case of Hibernate\index{Hibernate}).

Also the statements of several other query languages can be formatted in the same manner than SQL\index{SQL} statements.

\keeplisting{6}
The SQL \verb#SELECT#\index{SQL!SELECT} statement
\begin{lstlisting}[language=SQL]
SELECT person.id AS "id", person.name AS "lastName", person.surname AS "firstName", address.street AS "street", address.city AS "city"
    FROM contact.person, contact.address
    WHERE person.addressid = address.id
        AND person.profession like '%Teacher%';
\end{lstlisting}

\keeplisting{8}
would be embedded into a Java source file like this:
\begin{lstlisting}
final var sqlSelect = """
````SELECT person.id AS "id", person.name AS "lastName", person.surname AS "firstName", address.street AS "street", address.city AS "city"
````    FROM contact.person, contact.address
````    WHERE person.addressid = address.id
````        AND person.profession like '%Teacher%';\
````""";
\end{lstlisting}

Basically, the keywords are written all uppercase (\verb#SELECT#\index{SQL!SELECT}, \verb#FROM#\index{SQL!FROM},~…), while the arguments are all lowercase (\verb#person.name#, \verb#contact.address#,~…). Alias\index{SQL!alias} names are \textit{always} surrounded by double quotes even they do not contain blanks or are SQL\index{SQL} keywords\footnote{Some RDBMS\index{RDBMS} products do not allow double quotes for this purpose, they require square brackets.}.

The indendation\index{indentation} (four spaces)\index{indendation!character}\index{indentation!size} is per level: \verb#FROM#\index{SQL!FROM} and \verb#WHERE#\index{SQL!WHERE} are one level down from \verb#SELECT#\index{SQL!SELECT}, \verb#AND#\index{SQL!AND} is one level down from \verb#WHERE#\index{SQL!WHERE}.

%==============================================================================
% End of file!!
%==============================================================================

\subsection{Formatting HTML inside Java}\label{sec:FormattingHTMLInsideJava}
Embedded HTML\index{HTML} snippets are usually used to format the output of a program or an application, not only in the context of web applications.

Below some recommendations that were translated from \autocite{selfHTML:GuterHTMLStil}:
\begin{itemize}
	\item{Your HTML markup should be error-free (valid) so that it can be read by browsers and other parsers, such as screen readers. HTML\index{HTML} can only be validated if you include a \verb#doctype# declaration\footnote{This is only relevant if you embedd a complete HTML document into your Java code, and not just HTML snippets.}.} 
	\item{Virtues borrowed from XML\index{XML}: 
		\begin{itemize}
			\item{Well-formed syntax improves readability and enables processing as XHTML\index{XHTML}.}
			\item{All element names, attribute names, and their values are written in lowercase. Attribute values are always enclosed in doubel quotation marks.}
			\item{Use \bdq{meaningful} class names that describe function rather than visual styling.} 
			\item{Elements that do not strictly require a closing tag in HTML (e.g., \verb#p#, \verb#th#, \verb#td#, \verb#dt#, \verb#dd#, \verb#li#) should always be written with a closing tag.}
			\item{Use indentation and blank lines to structure your code clearly.}
		\end{itemize}}
\end{itemize}
\textbf{Note:} Since XHTML\index{XHTML} is an XML\index{XML} dialect, the points listed above are mandatory in XHTML\index{XHTML}, whereas in HTML\index{HTML} they are \bdq{merely} a matter of good practice.

\keeplisting{19}
This short HTML\index{HTML} document
\begin{lstlisting}[language=HTML]
<!DOCTYPE html>
<html>
    <head>
        <meta charset="UTF-8" />
        <title>Title</title>
    </head>
    <body>
        <h1>Headline</h2>
        <p>Text …</p>
        <p>More text …</p>
        <ol>
            <li>Entry 1</li>
            <li>Entry 2</li>
            <li>Entry 3</li>
        </ol>
        <p>Final text!</p>
    </body>
</html>
\end{lstlisting}

\keeplisting{21}
would be embedded like this:
\begin{lstlisting}
final var htmlDocumentString = """
````<!DOCTYPE html>
````<html>
````    <head>
````        <meta charset="UTF-8" />
````        <title>Title</title>
````    </head>
````    <body>
````        <h1>Headline</h2>
````        <p>Text …</p>
````        <p>More text …</p>
````        <ol>
````            <li>Entry 1</li>
````            <li>Entry 2</li>
````            <li>Entry 3</li>
````        </ol>
````        <p>Final text!</p>
````    </body>
````</html>
````""";
\end{lstlisting}

\keeplisting{11}
If you have to embed an HTML fragment – most probably, this would be a \verb#<p># element – this would look like this:
\begin{lstlisting}
final var htmlFragment = """
````<p>Text …</p>
````<p>More text …</p>
````    <ol>
````        <li>Entry 1</li>
````        <li>Entry 2</li>
````        <li>Entry 3</li>
````    </ol>
````<p>Final text!</p>
````""";
\end{lstlisting}

\keeplisting{4}
Or for an \verb#<li># element:
\begin{lstlisting}
final var htmlFragment = """
````<li>Entry 1</li>
````""";
\end{lstlisting}

This means that the top level element will not be indented. If you need to compose an HTML\index{HTML} document from the snippets, jyou can correct the indentation through calling \lstinline|String::indent|\index{java.lang!String!indent()} \autocite{ORACLE_DOC_STRING:indent}. This is explained in more detail in chapter \tqfullvref{sec:TextBlocks}.

%==============================================================================
% End of file!!
%==============================================================================


\subsection{Formatting XML inside Java}\label{sec:FormattingXMLInsideJava}
These formatting rules are valid for plain XML\index{XML}, as well as for various XML\index{XML} dialects, like SVG\index{SVG}, DocBook\index{DocBook} or XSLT\index{XSLT}, but not for XHTML\index{XHTML} – that is covered by \tqfullref{sec:FormattingHTMLInsideJava}, above.

The formatting of XML\index{XML} is relatively straight forward: the top level element will not be indented, each enclosed element will be indented by one level (or four spaces). The first attribute follows immediately the tag, further attributes are aligned with the first one. Attribute values are enclosed in double quotes. Closing tags are aligned with the opening tags, and the closing \bdq{\lstinline| />|} for an empty element is on the same line as the last attribute (or the opening tag, if there are no attributes).

\keeplisting{12}
This would look like this:
\begin{lstlisting}[language=XML]
<element attribute = "value">
    <inner attribute1 = "value"
           attribute2 = "value"
        <empty />
        <other attribute1 = "value"
               attribute2 = "value" />
    </inner>
</element>
<element attribute = "value">
    <inner attribute1 = "value" />
</element>
\end{lstlisting}

\keeplisting{14}
Embedded into Java code, the result would be this:
\begin{lstlisting}
final var xmlSnippet = """
````<element attribute = "value">
````    <inner attribute1 = "value"
````           attribute2 = "value"
````        <empty />
````        <other attribute1 = "value"
````               attribute2 = "value" />
````    </inner>
````</element>
````<element attribute = "value">
````    <inner attribute1 = "value" />
````</element>
````""";
\end{lstlisting}

%==============================================================================
% End of file!!
%==============================================================================


%==============================================================================
% End of file!!
%==============================================================================


%==============================================================================
% End of file!!
%==============================================================================


%==============================================================================
% End of file!!
%==============================================================================
